\section{Lists from own-question conditional means}
\label{ana:private-section}

The restriction lemma retains variance in conditional means that see
only one position's own questions and weights. We use this variance
to produce short lists of local labels. Set
\begin{equation}
 \beta=\frac\nu8,\qquad \ell=\frac2\beta,\qquad h=2\ell.
 \label{src:list-parameters}
\end{equation}
Since \(\nu=1/(2^7m^3)\), both \(\ell\) and \(h\) are
positive integers depending only on \(m\). Assume from now on that
the number \(d\) of graph positions also satisfies
\begin{equation}
 h<d,\qquad 16h+32h^2\leq\frac{\nu d}{2}.
 \label{src:restriction-budget}
\end{equation}
The joint parameter choice in \Cref{sec:conclusion} will ensure
these conditions. Applying \Cref{res:private-variance} to the
function in \Cref{ana:function} supplies a fixed set
\(J\subset[d]\) of size \(h\) and fixed data \(\boldsymbol a\)
consisting of the weights outside \(J\) and every seed outside
\(\mathcal E_J=\binom J2\). Write
\(f^*(\boldsymbol c_J,\boldsymbol s_J)\) for \(f\) with these
values plugged in. In the rest of this section the probability law
is the product of independent uniform hidden seeds
\((c_e)_{e\in\mathcal E_J}\) and independent uniform weight blocks
\((s_j)_{j\in J}\) on \([-1,1]^r\). The choice of the fixed data
does not condition these remaining draws on any event of success.

For \(j\in J\), write
\(u_j=(q_{jk})_{k\in J\setminus\{j\}}\), in increasing order of
\(k\), for the tuple of questions received by \(j\) from its
incident hidden seeds. Together with \(\boldsymbol a\), it
specifies the whole query \(i_j\); denote the completed query by
\(i_j(u_j)\). The conditional averages supplied by the restriction
lemma satisfy
\begin{equation}
 \sum_{j\in J}\E_{u_j}\Var_{s_j}H_j(u_j,s_j)
       \geq\nu h/4.
 \label{ana:restricted-variance}
\end{equation}
Here \(H_j(u,s)\) is the mean of \(f^*\) given the own questions
\(u_j=u\) and the own weight block \(s_j=s\).
Our aim is to turn this variance into own-input lists of size at most
\(\ell\) whose maximum compatible hit count has expectation at
least two.

To differentiate these means, we use the explicit integral from
\Cref{sec:restriction}. For a tuple \(u\) in
\(\supp u_j\), let \(\kappa_j(u)\) be the conditional distribution of
all hidden seeds given \(u_j=u\). On an incident pair
\(\{j,k\}\), this is uniform on the occurrence indices producing
the specified question at \(j\); on a nonincident hidden pair it
is uniform on \([M]\). All these conditional seed coordinates are
independent. This follows from independence of the original seeds
and from the fact that each condition fixes a function of just one
seed. Thus
\begin{equation}
 H_j(u,s)=
  \E_{\boldsymbol c_J\sim\kappa_j(u),\,(s_k)_{k\in J\setminus\{j\}}}
            f^*(\boldsymbol c_J,\boldsymbol s_J)
               \big|_{s_j=s},
 \label{ana:own-mean}
\end{equation}
where the other weight blocks have their original independent
uniform laws. The seed law in this formula does not depend on
\(s\). We use zero as a fallback for \(H_j\) on unsupported own
question tuples. The fixed data throughout are \(\Phi,F,J\), and
\(\boldsymbol a\).

\subsection{One gradient rule on every tuple}

The increments at positions outside \(J\) are fixed: their queries
use no seed in \(\mathcal E_J\). Set
\[
 c_0=\sum_{k\notin J}z(i_k,s_k).
\]
For a fixed tuple of batch queries \(\boldsymbol i=(i_j)_{j\in J}\),
consider the function on the whole real weight space
\[
 \psi_{\boldsymbol i}(\boldsymbol s_J)
       =F\left(c_0+\sum_{j\in J}z(i_j,s_j)\right).
\]
For \(I\in\mathcal I\), let
\(v_I(\boldsymbol i)\in\R^{J\times[r]}\) have entry one at
\((j,\operatorname{slot}_{i_j}(\lambda))\) if
\((i_j,\lambda)\in I\), and zero elsewhere. It has at most one
nonzero entry in each block \(j\).
We will represent the gradient using these incidence vectors. At
exceptional weights where a derivative need not exist, the same
representation will hold for a prescribed value of the gradient rule.

\begin{lemma}[A gradient rule in the compatibility hull]
\label{ana:gradient}
There is a deterministic measurable rule
\(g=(g_j)_{j\in J}\in\R^{J\times[r]}\), defined at every batch
query and real-weight tuple, that equals
\(\nabla\psi_{\boldsymbol i}\) outside a finite union of proper
affine hyperplanes for each fixed query tuple and fixed frozen
data. At every tuple,
\begin{equation}
 g\in\conv\{\pm v_I(\boldsymbol i):I\in\mathcal I\}.
 \label{ana:hull}
\end{equation}
Consequently,
\begin{equation}
 \|g_j\|_1\leq1,
 \qquad
 g_{j,t'}=0\quad\hbox{if }t'\notin
                   \operatorname{slot}_{i_j}(\Lambda_{i_j}).
 \label{ana:block-bound}
\end{equation}
The rule uses the fixed formula for \(F\), the frozen increments,
and the supplied batch queries and weights. It does not inspect
which hidden seed occurrences produced those queries.
\end{lemma}

\begin{proof}
Let \(\ind_I\in\R^{\mathcal P}\) be the indicator vector of
\(I\). For \(\zeta\in A_*\cup(-A_*)\),
\(I\in\mathcal I\), and \(\varepsilon\in\{-1,1\}\), form the
finite collection of affine expressions
\begin{equation}
 L_{\zeta,I,\varepsilon}(\boldsymbol s_J)
    =\varepsilon\langle\ind_I,c_0-\zeta\rangle
       +\varepsilon\langle v_I(\boldsymbol i),\boldsymbol s_J\rangle.
 \label{ana:affine-candidates}
\end{equation}
For either \(\mathcal C=A_*\) or \(\mathcal C=-A_*\),
\[
 \dist_*\left(c_0+\sum_{j\in J}z(i_j,s_j),\mathcal C\right)
   =\min_{\zeta\in\mathcal C}
       \max_{I\in\mathcal I,\,\varepsilon\in\{-1,1\}}
                    L_{\zeta,I,\varepsilon}(\boldsymbol s_J).
\]
In particular, \(\psi_{\boldsymbol i}\) is continuous and piecewise
affine, with finitely many pieces.

Compare all pairs of candidates in \eqref{ana:affine-candidates}.
If two candidates with different slope vectors in the unfrozen
weights tie in value, set \(g=0\). Otherwise take the true gradient
of the finite formula for \(\psi_{\boldsymbol i}\). This latter
choice is well-defined. Two candidates with the same slope either
never tie or are identical affine functions. All other comparisons
are strict at such a point and persist in a neighborhood; each
minimum and maximum there selects one affine function up to
identical copies. The resulting formula is affine in that
neighborhood. Fixing a deterministic order on candidates resolves
any choice between identical copies without changing its gradient.

Each discarded equality has a nonzero normal in the unfrozen
coordinates, hence defines a proper affine hyperplane. No choice of
frozen weights changes that fact, since frozen weights change the
offsets but not the slopes. Identical slopes caused by repeated
endpoint questions or other coincidences are allowed. Finite
comparisons prove measurability, and the rule is defined even at
query tuples that cannot arise from the hidden seed law. Its formula
uses the queries themselves and no additional seed information.

At a point not discarded, choose an active affine expression in
each distance formula. Write their slopes as
\(\varepsilon_-v_{I_-}(\boldsymbol i)\) for the distance to
\(-A_*\) and \(\varepsilon_+v_{I_+}(\boldsymbol i)\) for the
distance to \(A_*\). Since \(F\) is half the difference of
these distances,
\[
 g=\frac12\bigl(\varepsilon_-v_{I_-}(\boldsymbol i)
                  -\varepsilon_+v_{I_+}(\boldsymbol i)\bigr).
\]
Both vectors in this average are signed compatible-incidence
vectors, proving \eqref{ana:hull}. At discarded points the chosen
value zero also belongs to the symmetric hull. Every signed
incidence vector has block \(\ell^1\) norm at most one and zero
entries in the unused slots, so convexity proves
\eqref{ana:block-bound} at every tuple.
\end{proof}

\subsection{Own-question conditional gradients}

For supported \(u\) and any \(s\in\R^r\), define
\begin{equation}
 \bar g_j(u,s)=
   \E_{\boldsymbol c_J\sim\kappa_j(u),\,(s_k)_{k\in J\setminus\{j\}}}
               g_j(\boldsymbol c_J,\boldsymbol s_J)\big|_{s_j=s}.
 \label{ana:conditional-gradient}
\end{equation}
Here the notation for \(g_j\) means that the seed tuple is used
only to form the queries supplied to the rule in
\Cref{ana:gradient}. Set \(\bar g_j(u,s)=0\) for unsupported
own tuples. For supported \(u\), its query \(i_j(u)\) is fixed
throughout the integral. Convexity and \eqref{ana:block-bound} give,
at every own input,
\begin{equation}
 \|\bar g_j(u,s)\|_1\leq1,
 \qquad
 \bar g_{j,t'}(u,s)=0\quad\hbox{on unused slots of }i_j(u).
 \label{ana:mean-block-bound}
\end{equation}

\begin{lemma}[Conditional cube variance]
\label{ana:cube}
For each \(j\in J\) and supported own tuple \(u\), the function
\(H_j(u,\cdot)\) is \(1\)-Lipschitz with respect to
\(\|\cdot\|_\infty\). For almost every own
\(s\in[-1,1]^r\),
\begin{equation}
 \partial_{s^{(t')}}H_j(u,s)=\bar g_{j,t'}(u,s),
 \qquad t'\in[r].
 \label{ana:conditional-derivative}
\end{equation}
Its variance satisfies the dimension-independent bound
\begin{equation}
 \Var_s H_j(u,s)
       \leq2\E_s\|\bar g_j(u,s)\|_2^2.
 \label{ana:cube-variance}
\end{equation}
Consequently, under the fixed restricted law,
\begin{equation}
 \sum_{j\in J}\E\|\bar g_j(u_j,s_j)\|_2^2
      \geq\nu h/8=\beta h.
 \label{ana:energy}
\end{equation}
\end{lemma}

\begin{proof}
Replacing the own block \(s\) by \(s'\), with a seed configuration
and other weights fixed, changes the integrand in
\eqref{ana:own-mean} by at most \(\|s-s'\|_\infty\). The same
bound holds after averaging. In particular, every scalar-coordinate
difference quotient of the integrand is bounded in absolute value
by one.

For each fixed hidden seed tuple, the full gradient identification
in \Cref{ana:gradient} holds outside a null set in all batch weights.
Fubini implies that for almost every own \(s\) it holds for almost
all other weight blocks. The conditional seed law \(\kappa_j(u)\) has
finite support and does not depend on \(s\), so the exceptional sets
may be united over its support and over the finitely many weight
coordinates. At the remaining own values, the coordinate difference
quotients converge almost surely to \(g_{j,t'}\). Dominated
convergence gives \eqref{ana:conditional-derivative}; the boundary
of the cube has measure zero.

Every coordinate-line restriction of \(H_j(u,\cdot)\) is
Lipschitz and hence absolutely continuous. By Fubini applied to
\eqref{ana:conditional-derivative}, for almost every fixing of the
other own coordinates its derivative agrees almost everywhere with
the corresponding entry of \(\bar g_j\). The fundamental theorem
of calculus therefore expresses differences between any two values
on such a line as the integral of that entry.

Let \(w_0\) be the restriction of \(H_j(u,\cdot)\) to one of these
lines. For \(x,y\in[-1,1]\), Cauchy--Schwarz gives
\[
 |w_0(x)-w_0(y)|^2
    \leq |x-y|\int_{\min(x,y)}^{\max(x,y)}|w_0'(z)|^2\,dz
    \leq2\int_{-1}^1|w_0'(z)|^2\,dz.
\]
For independent uniform \(X,Y\in[-1,1]\), this implies
\[
 \Var(w_0(X))=\tfrac12\E|w_0(X)-w_0(Y)|^2
      \leq\int_{-1}^1|w_0'(z)|^2\,dz
      =2\E|w_0'(X)|^2.
\]
The product decomposition of Efron and
Stein~\cite[Section~2]{EfronStein1981} bounds variance by the sum
of the averaged one-coordinate conditional variances. To see this
consequence directly,
decompose the identity at each scalar coordinate into its averaging
projection and its orthogonal complement. The resulting orthogonal
components are indexed by subsets of coordinates. Total variance
counts every nonconstant squared component once, while the sum of
one-coordinate conditional variances counts it once for each member
of its nonempty index subset. Applying the line estimate,
integrating over the other coordinates, and summing proves
\eqref{ana:cube-variance}. Averaging that bound over the own tuples
and summing over \(j\), then using
\eqref{ana:restricted-variance}, proves \eqref{ana:energy}.
\end{proof}

A particular own weight value can make an entire slice of the other
weights exceptional. The derivative identity is only an
almost-everywhere assertion, and \eqref{ana:energy} is integrated
over the restricted law. In contrast,
\eqref{ana:hull}, \eqref{ana:block-bound}, and
\eqref{ana:mean-block-bound} hold at every input of their respective
rules, including exceptional weights.

\subsection{Thresholding the conditional gradients}

The conditional means now retain total squared gradient energy at
least \(\beta h\), and each block has \(\ell^1\) norm at most one.
Keeping coordinates above a fixed threshold therefore gives short
lists without an alphabet-size factor. The remaining task is to
transfer their mass back to the single full gradient, whose compatible
incidence representation can certify simultaneous hits.

For supported \(u\), define the slot-ordered list
\begin{equation}
 \mathcal L_j(u,s)=
  \left\{\lambda\in\Lambda_{i_j(u)}:
    \left|\bar g_{j,\operatorname{slot}_{i_j(u)}(\lambda)}(u,s)\right|
       \geq\frac\beta2\right\}.
 \label{ana:list-definition}
\end{equation}
For an unsupported own question tuple, let the list be empty.
At an actual draw write \(L_j=\mathcal L_j(u_j,s_j)\), and define
\begin{equation}
 Z=\max_{I\in\mathcal I}
       \#\{j\in J:\exists\lambda\in L_j\text{ with }(i_j,\lambda)\in I\}.
 \label{ana:hit-count}
\end{equation}

\begin{proposition}[Local lists and compatible hits]
\label{ana:lists}
The functions \(\mathcal L_j\) are deterministic total measurable maps
returning at most \(\ell\) valid labels. Once
\(\Phi,F,J,\boldsymbol a\) and the conditional kernels above are
fixed, each list uses only its own questions \(u_j\) and weights
\(s_j\). In particular it has no additional access to the actual
hidden seed occurrences or to any other actual batch input.
Under the fixed restricted product law,
\begin{equation}
 \E Z\geq\frac{\beta h}{2}=2.
 \label{ana:list-hits}
\end{equation}
\end{proposition}

\begin{proof}
The mean block bound gives
\[
 \frac\beta2|\mathcal L_j(u,s)|\leq\|\bar g_j(u,s)\|_1\leq1,
 \qquad |\mathcal L_j(u,s)|\leq2/\beta=\ell.
\]
It also gives zero dummy entries, so thresholding returns only valid
labels. Finite threshold comparisons of measurable functions, with
an empty-list fallback, prove totality and measurability.

We first transfer the retained mass of the own conditional means
to the full gradient rule \(g\), and then apply its pointwise hull bound.
For an actual own input, let \(T_j\subset[r]\) be the set of slots
of \(L_j\), and write \(\mathcal G_j\) for the information
\((u_j,s_j)\). The conditional mean in
\eqref{ana:conditional-gradient} is a version of
\(\E[g_j\mid\mathcal G_j]\), defined by the displayed kernel even
at own values of probability zero. For slots outside \(T_j\),
\eqref{ana:mean-block-bound} and the threshold give
\[
 \sum_{t'\notin T_j}|\bar g_{j,t'}|^2
       \leq\frac\beta2\sum_{t'\notin T_j}|\bar g_{j,t'}|
       \leq\frac\beta2.
\]
For the list slots, each mean entry has magnitude at most one, and
conditional Jensen gives
\[
 \sum_{t'\in T_j}|\bar g_{j,t'}|^2
   \leq\sum_{t'\in T_j}|\bar g_{j,t'}|
   \leq\E\left[\sum_{t'\in T_j}|g_{j,t'}|
                         \,\middle|\,\mathcal G_j\right].
\]
The last inequality uses the fact that \(T_j\) is determined by
\(\mathcal G_j\). Taking expectations, summing, and applying
\eqref{ana:energy} therefore yields
\begin{equation}
 \E\sum_{j\in J}\sum_{\lambda\in L_j}
        |g_{j,\operatorname{slot}_{i_j}(\lambda)}|
       \geq\beta h-\frac\beta2h=\frac{\beta h}{2}.
 \label{ana:retained-mass}
\end{equation}

For arbitrary fixed queries and arbitrary lists of valid labels,
the sum of absolute values over their slot sets is a convex
function of the full vector in \(\R^{J\times[r]}\). On a signed
generator \(\pm v_I(\boldsymbol i)\) it equals the number of
positions hit by that selection \(I\), since each block contains
at most one selected slot. Its value on every point of the hull
in \eqref{ana:hull} is consequently at most the maximum such hit
count. At every realized tuple, apply this deterministic statement
to the realized lists and the full vector \(g\):
\begin{equation}
 \sum_{j\in J}\sum_{\lambda\in L_j}
        |g_{j,\operatorname{slot}_{i_j}(\lambda)}|\leq Z.
 \label{ana:pointwise-hits}
\end{equation}
No independence between the lists and the gradient, or between
positions, is required. Combining \eqref{ana:retained-mass} and
\eqref{ana:pointwise-hits} proves \eqref{ana:list-hits}, with
\(\beta h/2=2\) from \eqref{src:list-parameters}.

Finally, \eqref{ana:own-mean} and
\eqref{ana:conditional-gradient} integrate over hidden seeds
against the fixed law \(\kappa_j(u)\), and over other weights against
their fixed cube laws. Their actual realized values are not inputs
to the resulting mean function. Thus \eqref{ana:list-definition}
has exactly the stated own-input dependence. No efficiency of this
rule is needed for the soundness argument.
\end{proof}

The kernels and list maps in this section remain fixed throughout
the subsequent argument. Conditioning later on other hidden seeds
or on all weights changes the distribution of the inputs to these
maps; it does not recompute their conditional averages. Locality,
validity, and list-size bounds continue to hold for every fixed
weight choice. The lower bound \eqref{ana:list-hits} concerns the
original restricted product law.
